% ============================================
% Supplementary Figures
% ============================================
% Reordered to match sequential appearance in main text
% Note: Some figures moved to Extended Data (ED Figs 1-8)
\newpage
\subsection*{Supplementary Figures}

% ============================================
% Figures 1-2: Strain characterization (Section 3.1)
% ============================================

% Supple Fig 1 - Taxonomy color map
\begin{figure}[H]
\centering
\includegraphics[width=0.9\textwidth]{supplementary_figs/taxonomy_color_map.pdf}
\caption{\textbf{Supplementary Fig.~1.} Taxonomic classification of bacterial isolates. Table showing the 54 bacterial isolates used in the experiment, of which 4 pairs share identical ASV sequences. Each row shows the ASV identifier and its full taxonomic classification (Kingdom, Phylum, Class, Order, Family, Genus). Colors match those used in pie charts and composition figures throughout the manuscript.}
\label{fig:suppl_s1}
\end{figure}

% Supple Fig 2 - Phylogeny tree
\begin{figure}[H]
\centering
\includegraphics[width=0.9\textwidth]{supplementary_figs/Fig_1-3_phylogeny_tree.pdf}
\caption{\textbf{Supplementary Fig.~2.} Phylogenetic tree of bacterial isolates. Maximum-likelihood phylogenetic tree based on full-length 16S rRNA gene sequences of the 54 bacterial isolates used in this study. Tree was constructed using EMBL-EBI tools\citep{Madeira2022}. The isolates span three phyla (Proteobacteria, Firmicutes, Bacteroidota) and 29 families, representing broad phylogenetic diversity.}
\label{fig:suppl_s2}
\end{figure}

% Note: Time series in Base media moved to Extended Data Fig. 6
% Note: Robustness of outcome classification moved to Extended Data Fig. 7
% Note: Skewness null model comparison moved to Extended Data Fig. 1

% ============================================
% Figures 3-5: Simulation robustness (Section 3.3)
% ============================================

% Supple Fig 3 - Simulation: growth-rate heterogeneity
\begin{figure}[H]
\centering
\begin{subfigure}[b]{0.48\textwidth}
    \includegraphics[width=\textwidth]{supplementary_figs/Fig_phase_diagram_ablation_growth_std01.pdf}
    \caption{Growth rate: mean 1, standard deviation 0.1}
\end{subfigure}
\hfill
\begin{subfigure}[b]{0.48\textwidth}
    \includegraphics[width=\textwidth]{supplementary_figs/Fig_phase_diagram_ablation_growth_std02.pdf}
    \caption{Growth rate: mean 1, standard deviation 0.2}
\end{subfigure}
\caption{\textbf{Supplementary Fig.~3.} Phase diagrams with growth-rate heterogeneity. Coalescence outcomes when species have heterogeneous intrinsic growth rates sampled from normal distributions with mean 1 and standard deviations of \textbf{(a)} 0.1 and \textbf{(b)} 0.2. The qualitative transition from Mixture-dominated to Dominance/Restructuring outcomes with increasing interaction strength is robust to growth-rate variation.}
\label{fig:suppl_s3}
\end{figure}

% Supple Fig 4 - Simulation: carrying-capacity variation
\begin{figure}[H]
\centering
\begin{subfigure}[b]{0.48\textwidth}
    \includegraphics[width=\textwidth]{supplementary_figs/Fig_phase_diagram_ablation_k_std01.pdf}
    \caption{Carrying capacity: mean 1, standard deviation 0.1}
\end{subfigure}
\hfill
\begin{subfigure}[b]{0.48\textwidth}
    \includegraphics[width=\textwidth]{supplementary_figs/Fig_phase_diagram_ablation_k_std02.pdf}
    \caption{Carrying capacity: mean 1, standard deviation 0.2}
\end{subfigure}
\caption{\textbf{Supplementary Fig.~4.} Phase diagrams with carrying-capacity variation. Coalescence outcomes when species have heterogeneous carrying capacities sampled from normal distributions with mean 1 and standard deviations of \textbf{(a)} 0.1 and \textbf{(b)} 0.2. The qualitative transition from Mixture-dominated to Dominance/Restructuring outcomes with increasing interaction strength is robust to carrying-capacity variation.}
\label{fig:suppl_s4}
\end{figure}

% Supple Fig 5 - Simulation: alternative alpha_ij distributions
\begin{figure}[H]
\centering
\begin{subfigure}[b]{0.48\textwidth}
    \includegraphics[width=\textwidth]{supplementary_figs/Fig_phase_diagram_ablation_gaussian.pdf}
    \caption{Gaussian distribution (mean $\mu$, matched variance)}
\end{subfigure}
\hfill
\begin{subfigure}[b]{0.48\textwidth}
    \includegraphics[width=\textwidth]{supplementary_figs/Fig_phase_diagram_ablation_gamma.pdf}
    \caption{Gamma distribution (mean $\mu$, matched variance)}
\end{subfigure}
\caption{\textbf{Supplementary Fig.~5.} Phase diagrams for alternative interaction coefficient distributions. Both distributions have the same mean and variance but different shapes. \textbf{(a)} Gaussian distribution (truncated at 0 to ensure non-negative coefficients). \textbf{(b)} Gamma distribution (inherently non-negative). The qualitative transition from Mixture-dominated to Dominance/Restructuring outcomes with increasing mean interaction strength $\mu$ is robust to the choice of distribution, confirming that the observed patterns depend primarily on interaction strength rather than distributional details.}
\label{fig:suppl_s5}
\end{figure}

% Note: Species number ablation moved to Extended Data Fig. 4
% Note: Pairwise selection correlation (simulation) moved to Extended Data Fig. 2

% ============================================
% Figures 6-8: Pairwise invasion matrices (Section 3.4)
% ============================================

% Supple Fig 6 - Nutr- pairwise invasion
\begin{figure}[H]
\centering
\includegraphics[width=0.9\textwidth]{supplementary_figs/Pairwise_Matrix_Dynamics_LN_improved.pdf}
\caption{\textbf{Supplementary Fig.~6.} Pairwise invasion outcomes in Nutr$-$ medium. Matrix showing competitive outcomes among the 12 most abundant isolates. Each cell indicates whether the invader (initially 5\%) successfully established or was excluded after seven growth-dilution cycles. Low failed-invasion frequency indicates weak competitive exclusion.}
\label{fig:suppl_s6}
\end{figure}

% Supple Fig 7 - Base pairwise invasion
\begin{figure}[H]
\centering
\includegraphics[width=0.9\textwidth]{supplementary_figs/Pairwise_Matrix_Dynamics_MN_improved.pdf}
\caption{\textbf{Supplementary Fig.~7.} Pairwise invasion outcomes in Base medium. Matrix showing competitive outcomes among the 12 most abundant isolates. Each cell indicates whether the invader (initially 5\%) successfully established or was excluded after seven growth-dilution cycles. Intermediate failed-invasion frequency indicates moderate competitive exclusion.}
\label{fig:suppl_s7}
\end{figure}

% Supple Fig 8 - Nutr+ pairwise invasion
\begin{figure}[H]
\centering
\includegraphics[width=0.9\textwidth]{supplementary_figs/Pairwise_Matrix_Dynamics_HN_improved.pdf}
\caption{\textbf{Supplementary Fig.~8.} Pairwise invasion outcomes in Nutr$+$ medium. Matrix showing competitive outcomes among the 12 most abundant isolates. Each cell indicates whether the invader (initially 5\%) successfully established or was excluded after seven growth-dilution cycles. High failed-invasion frequency indicates strong competitive exclusion.}
\label{fig:suppl_s8}
\end{figure}

% Note: Pairwise selection correlation (experimental) moved to Extended Data Fig. 3

% ============================================
% Figures 9-10: pH mechanism (Section 3.5)
% ============================================

% Supple Fig 9 - Species-pH correlation
\begin{figure}[H]
\centering
\includegraphics[width=0.9\textwidth]{supplementary_figs/two_panel_ph_asv_figure.pdf}
\caption{\textbf{Supplementary Fig.~9.} Monoculture pH modification by bacterial isolates. Distribution of monoculture pH after 15h growth for all ASVs, showing the range of pH modification abilities relative to the initial pH of 6.5. Strong acidifiers (monoculture pH $<$ 5) include ASV 12 (\textit{Serratia}), ASV 3 (\textit{Enterobacter}), and ASV 8 (\textit{Citrobacter}). Species that maintain or slightly increase pH above initial (monoculture pH $>$ 6.5) include ASV 11 (\textit{Stenotrophomonas}) and ASV 9 (\textit{Pseudomonas}). Bottom legend shows taxa identities for the labeled ASVs.}
\label{fig:suppl_s9}
\end{figure}

% Supple Fig 10 - ASV abundance vs parent community pH
\begin{figure}[H]
\centering
\includegraphics[width=0.95\textwidth]{supplementary_figs/Fig_S23_ASV_vs_pH_combined.pdf}
\caption{\textbf{Supplementary Fig.~10.} Dominant species class (acidifier vs.\ alkalizer) determines parental community pH. Scatter plots showing the relationship between dominant ASV relative abundance in parental communities and the resulting community pH at stabilization. Top row: alkalizers including ASV 9 (\textit{Pseudomonas}) and ASV 11 (\textit{Stenotrophomonas}), whose dominance is associated with higher community pH. Bottom row: acidifiers including ASV 3 (\textit{Enterobacter}), ASV 8 (\textit{Citrobacter}), and ASV 12 (\textit{Serratia}), whose dominance is associated with lower community pH. For each ASV, left panel shows Base medium and right panel shows Nutr$+$ medium. Regression lines and correlation coefficients shown; asterisks indicate $p < 0.05$.}
\label{fig:suppl_s10}
\end{figure}

% Note: pH difference vs coalescence outcome moved to Extended Data Fig. 8

% ============================================
% Figures 11-13: Rank-Abundance Curves
% ============================================

% Supple Fig 11 - Rank-Abundance Curves - Base medium
\begin{figure}[H]
\centering
\includegraphics[width=0.85\textwidth]{supplementary_figs/rank_abundance_parental_vs_coalesced_Base.pdf}
\caption{\textbf{Supplementary Fig.~11.} Rank-abundance curves comparing parental and coalesced communities in Base medium. \textbf{(a)} Parental (sub)communities before coalescence. \textbf{(b)} Coalesced communities after coalescence. Each line represents one community's rank-abundance curve. Gini coefficients quantify abundance inequality (0 = perfectly even, 1 = one species dominates). Parental: Gini = 0.62 $\pm$ 0.19 (n=59); Coalesced: Gini = 0.64 $\pm$ 0.15 (n=94). Both parental and coalesced communities exhibit substantial abundance skewness with similar Gini values.}
\label{fig:suppl_s11}
\end{figure}

% Supple Fig 12 - Rank-Abundance Curves - Nutr- medium
\begin{figure}[H]
\centering
\includegraphics[width=0.85\textwidth]{supplementary_figs/rank_abundance_parental_vs_coalesced_Nutr-.pdf}
\caption{\textbf{Supplementary Fig.~12.} Rank-abundance curves comparing parental and coalesced communities in Nutr$-$ medium. \textbf{(a)} Parental (sub)communities before coalescence. \textbf{(b)} Coalesced communities after coalescence. Each line represents one community's rank-abundance curve. Parental: Gini = 0.63 $\pm$ 0.09 (n=60); Coalesced: Gini = 0.61 $\pm$ 0.12 (n=92).}
\label{fig:suppl_s12}
\end{figure}

% Supple Fig 13 - Rank-Abundance Curves - Nutr+ medium
\begin{figure}[H]
\centering
\includegraphics[width=0.85\textwidth]{supplementary_figs/rank_abundance_parental_vs_coalesced_Nutr+.pdf}
\caption{\textbf{Supplementary Fig.~13.} Rank-abundance curves comparing parental and coalesced communities in Nutr$+$ medium. \textbf{(a)} Parental (sub)communities before coalescence. \textbf{(b)} Coalesced communities after coalescence. Each line represents one community's rank-abundance curve. Parental: Gini = 0.64 $\pm$ 0.17 (n=60); Coalesced: Gini = 0.66 $\pm$ 0.14 (n=94).}
\label{fig:suppl_s13}
\end{figure}

% ============================================
% Figures 14-16: Coalescence Outcome Matrices
% ============================================

% Supple Fig 14 - Coalescence Matrix - Nutr-
\begin{figure}[H]
\centering
\begin{subfigure}[b]{0.72\textwidth}
    \includegraphics[width=\textwidth]{supplementary_figs/PieCharts/CoalescenceMatrices/coalescence_matrix_Nutr-_s6.pdf}
    \caption{Initial richness: 6 species}
\end{subfigure}

\vspace{0.5cm}

\begin{subfigure}[b]{0.72\textwidth}
    \includegraphics[width=\textwidth]{supplementary_figs/PieCharts/CoalescenceMatrices/coalescence_matrix_Nutr-_s12.pdf}
    \caption{Initial richness: 12 species}
\end{subfigure}

\vspace{0.5cm}

\begin{subfigure}[b]{0.72\textwidth}
    \includegraphics[width=\textwidth]{supplementary_figs/PieCharts/CoalescenceMatrices/coalescence_matrix_Nutr-_s24.pdf}
    \caption{Initial richness: 24 species}
\end{subfigure}
\caption{\textbf{Supplementary Fig.~14.} Coalescence outcome matrices in Nutr$-$ medium. Each matrix shows pairwise coalescence outcomes for all parental community combinations at different initial richness levels. Pie charts display the species composition of coalesced communities; colors match parental origins. Under weak interactions, outcomes show more balanced mixing.}
\label{fig:suppl_s14}
\end{figure}

% Supple Fig 15 - Coalescence Matrix - Base
\begin{figure}[H]
\centering
\begin{subfigure}[b]{0.72\textwidth}
    \includegraphics[width=\textwidth]{supplementary_figs/PieCharts/CoalescenceMatrices/coalescence_matrix_Base_s6.pdf}
    \caption{Initial richness: 6 species}
\end{subfigure}

\vspace{0.5cm}

\begin{subfigure}[b]{0.72\textwidth}
    \includegraphics[width=\textwidth]{supplementary_figs/PieCharts/CoalescenceMatrices/coalescence_matrix_Base_s12.pdf}
    \caption{Initial richness: 12 species}
\end{subfigure}

\vspace{0.5cm}

\begin{subfigure}[b]{0.72\textwidth}
    \includegraphics[width=\textwidth]{supplementary_figs/PieCharts/CoalescenceMatrices/coalescence_matrix_Base_s24.pdf}
    \caption{Initial richness: 24 species}
\end{subfigure}
\caption{\textbf{Supplementary Fig.~15.} Coalescence outcome matrices in Base medium. Each matrix shows pairwise coalescence outcomes for all parental community combinations at different initial richness levels. Pie charts display the species composition of coalesced communities; colors match parental origins. Moderate interactions produce frequent one-sided outcomes (Dominance).}
\label{fig:suppl_s15}
\end{figure}

% Supple Fig 16 - Coalescence Matrix - Nutr+
\begin{figure}[H]
\centering
\begin{subfigure}[b]{0.72\textwidth}
    \includegraphics[width=\textwidth]{supplementary_figs/PieCharts/CoalescenceMatrices/coalescence_matrix_Nutr+_s6.pdf}
    \caption{Initial richness: 6 species}
\end{subfigure}

\vspace{0.5cm}

\begin{subfigure}[b]{0.72\textwidth}
    \includegraphics[width=\textwidth]{supplementary_figs/PieCharts/CoalescenceMatrices/coalescence_matrix_Nutr+_s12.pdf}
    \caption{Initial richness: 12 species}
\end{subfigure}

\vspace{0.5cm}

\begin{subfigure}[b]{0.72\textwidth}
    \includegraphics[width=\textwidth]{supplementary_figs/PieCharts/CoalescenceMatrices/coalescence_matrix_Nutr+_s24.pdf}
    \caption{Initial richness: 24 species}
\end{subfigure}
\caption{\textbf{Supplementary Fig.~16.} Coalescence outcome matrices in Nutr$+$ medium. Each matrix shows pairwise coalescence outcomes for all parental community combinations at different initial richness levels. Pie charts display the species composition of coalesced communities; colors match parental origins. Strong interactions produce predominantly one-sided outcomes (Dominance).}
\label{fig:suppl_s16}
\end{figure}

% ============================================
% Figures 17-18: Additional Time Series
% ============================================

% Supple Fig 17 - Nutr- time series
\begin{figure}[H]
\centering
\begin{subfigure}[b]{0.45\textwidth}
    \includegraphics[width=\textwidth]{supplementary_figs/timeseries_merged/Fig_1-3_timeseries_Nutr-_ex1_merged.pdf}
    \caption{}
\end{subfigure}
\hfill
\begin{subfigure}[b]{0.45\textwidth}
    \includegraphics[width=\textwidth]{supplementary_figs/timeseries_merged/Fig_1-3_timeseries_Nutr-_ex2_merged.pdf}
    \caption{}
\end{subfigure}
\caption{\textbf{Supplementary Fig.~17.} Time series of community composition during coalescence in Nutr$-$ medium. Each panel \textbf{(a, b)} shows two parental subcommunities (left, stacked vertically) and their coalesced outcome (right) over serial dilution cycles. Stacked bar charts represent relative species abundances at each time point. X-axis: serial dilution cycles (0--7); Y-axis: relative abundance.}
\label{fig:suppl_s17}
\end{figure}

% Supple Fig 18 - Nutr+ time series
\begin{figure}[H]
\centering
\includegraphics[width=0.5\textwidth]{supplementary_figs/timeseries_merged/Fig_1-3_timeseries_Nutr+_ex1_merged.pdf}
\caption{\textbf{Supplementary Fig.~18.} Time series of community composition during coalescence in Nutr$+$ medium. The panel shows two parental subcommunities (left, stacked vertically) and their coalesced outcome (right) over serial dilution cycles. Stacked bar charts represent relative species abundances at each time point. X-axis: serial dilution cycles (0--7); Y-axis: relative abundance.}
\label{fig:suppl_s18}
\end{figure}

% ============================================
% Figures 19-20: Assembly Effect
% ============================================

% Supple Fig 19 - Assembly reduces interaction strength
\begin{figure}[H]
\centering
\includegraphics[width=0.6\textwidth]{supplementary_figs/Assembly_effect_scatter_combined.pdf}
\caption{\textbf{Supplementary Fig.~19.} Assembly reduces mean pairwise interaction strength. Mean interaction strength before assembly (red, inter-community) versus after assembly (blue, intra-community) across different interaction strength parameters $\mu$. Points show individual simulation runs; lines connect means. Post-assembly communities consistently show reduced interaction strength, indicating that assembly filters out strongly competing species.}
\label{fig:suppl_s19}
\end{figure}

% Note: Assembly effect simulation stacked bar moved to Extended Data Fig. 5

% Supple Fig 20 - Assembly effect experimental
\begin{figure}[H]
\centering
\includegraphics[width=0.7\textwidth]{supplementary_figs/Fig_S26_assembly_effect_experimental_stacked_bar.pdf}
\caption{\textbf{Supplementary Fig.~20.} Assembly history promotes Dominance in experiments. Stacked bar plots comparing outcome distributions between coalescence (top row) and direct assembly (bottom row) across three nutrient conditions (Nutr$-$, Base, Nutr$+$). Percentages and counts (n/total) are shown for each category. In Base and Nutr$+$ conditions, coalescence shows elevated Dominance rates compared to direct assembly, consistent with simulation predictions (Extended Data Fig.~5, Supplementary Fig.~19).}
\label{fig:suppl_s20}
\end{figure}

% ============================================
% Figures 21-22: Monoculture and Overlap
% ============================================

% Supple Fig 21 - Monoculture OD and growth rate characterization
\begin{figure}[H]
\centering
\includegraphics[width=0.9\textwidth]{supplementary_figs/monoculture_od_growth_histograms.pdf}
\caption{\textbf{Supplementary Fig.~21.} Phenotypic diversity of bacterial isolates in monoculture (Base medium). \textbf{(a)} Distribution of optical density (OD$_{600}$) reached by bacterial isolates during monoculture growth in Base medium ($n = 54$ isolates). \textbf{(b)} Distribution of exponential growth rates across isolates ($n = 45$ isolates with measurable growth rates out of 54 total). Growth rates were computed by log-linear regression of OD versus time during the exponential phase from full kinetic time-course measurements (475 cycles, $\sim$3~min intervals). The isolates exhibit broad variation in both growth yield and growth rate, reflecting phenotypic diversity that may contribute to varied coalescence outcomes in community mixing experiments.}
\label{fig:suppl_s21}
\end{figure}

% Supple Fig 22 - Overlap fraction histogram
\begin{figure}[H]
\centering
\includegraphics[width=0.95\textwidth]{supplementary_figs/overlap_fraction_histogram.pdf}
\caption{\textbf{Supplementary Fig.~22.} Fraction of coalesced ASVs present in both parental communities. Distribution of overlap fraction across coalescence events in three nutrient conditions. Overlap fraction is defined as the proportion of ASVs in the coalesced community that were present in both parental communities before mixing. Nutr$-$: mean = 0.10 $\pm$ 0.10 (n=90); Base: mean = 0.22 $\pm$ 0.26 (n=83); Nutr$+$: mean = 0.17 $\pm$ 0.25 (n=90). The relatively low overlap fractions across all conditions indicate that most surviving ASVs in coalesced communities originated from only one of the two parental communities, consistent with the prevalence of Dominance outcomes.}
\label{fig:suppl_s22}
\end{figure}

% ============================================
% Figures 23-26: Natural Community (Section 3.6)
% ============================================

% Supple Fig 23 - Natural rank-abundance Nutr-
\begin{figure}[H]
\centering
\includegraphics[width=0.85\textwidth]{supplementary_figs/rank_abundance_natural_Nutr-.pdf}
\caption{\textbf{Supplementary Fig.~23.} Rank-abundance curves for natural sample-derived communities in Nutr$-$ medium. \textbf{(a)} Parental (sub)communities before coalescence. \textbf{(b)} Coalesced communities after coalescence. Each line represents one community's rank-abundance curve. Gini coefficients quantify abundance inequality (0 = perfectly even, 1 = one species dominates). ASV richness indicates number of ASVs above 0.1\% relative abundance threshold. Horizontal dashed line indicates 0.1\% threshold. Natural communities exhibit higher ASV richness compared to synthetic communities (which had controlled initial richness of 6, 12, or 24 species), reflecting the complex species assemblages in environmental samples.}
\label{fig:suppl_s23}
\end{figure}

% Supple Fig 24 - Natural rank-abundance Base
\begin{figure}[H]
\centering
\includegraphics[width=0.85\textwidth]{supplementary_figs/rank_abundance_natural_Base.pdf}
\caption{\textbf{Supplementary Fig.~24.} Rank-abundance curves for natural sample-derived communities in Base medium (natural-community analog of Supplementary Fig.~12 for synthetic communities). \textbf{(a)} Parental (sub)communities before coalescence. \textbf{(b)} Coalesced communities after coalescence. Each line represents one community's rank-abundance curve. Gini coefficients quantify abundance inequality (0 = perfectly even, 1 = one species dominates). ASV richness indicates number of ASVs above 0.1\% relative abundance threshold. Horizontal dashed line indicates 0.1\% threshold.}
\label{fig:suppl_s24}
\end{figure}

% Supple Fig 25 - Natural rank-abundance Nutr+
\begin{figure}[H]
\centering
\includegraphics[width=0.85\textwidth]{supplementary_figs/rank_abundance_natural_Nutr+.pdf}
\caption{\textbf{Supplementary Fig.~25.} Rank-abundance curves for natural sample-derived communities in Nutr$+$ medium (natural-community analog of Supplementary Fig.~13 for synthetic communities). \textbf{(a)} Parental (sub)communities before coalescence. \textbf{(b)} Coalesced communities after coalescence. Each line represents one community's rank-abundance curve. Gini coefficients quantify abundance inequality (0 = perfectly even, 1 = one species dominates). ASV richness indicates number of ASVs above 0.1\% relative abundance threshold. Horizontal dashed line indicates 0.1\% threshold.}
\label{fig:suppl_s25}
\end{figure}

% Supple Fig 26 - Natural overlap fraction histogram
\begin{figure}[H]
\centering
\includegraphics[width=0.95\textwidth]{supplementary_figs/overlap_fraction_histogram_natural.pdf}
\caption{\textbf{Supplementary Fig.~26.} Fraction of coalesced ASVs present in both parental communities for natural sample-derived communities. Distribution of overlap fraction across coalescence events in three nutrient conditions. Overlap fraction is defined as the proportion of ASVs in the coalesced community that were present in both parental communities before mixing. Nutr$-$: mean = 0.13 $\pm$ 0.09 (n=30); Base: mean = 0.09 $\pm$ 0.13 (n=30); Nutr$+$: mean = 0.23 $\pm$ 0.15 (n=30). Similar to synthetic communities (Supplementary Fig.\ 22), the relatively low overlap fractions indicate that most surviving ASVs in coalesced communities originated from only one of the two parental communities, consistent with the prevalence of Dominance outcomes.}
\label{fig:suppl_s26}
\end{figure}
